\SourceLeaf{357}
\par\bigskip\Needspace{7\baselineskip}
\pdfbookmark[1]{Zur Theorie der biquadratischen Reste. V.}{nachlass-biquadratic-V}
\centerline{[V.]}\medskip
\renewcommand{\Article}[1]{\par\medskip\Needspace{4\baselineskip}\pdfbookmark[2]{Art. #1}{nachlass-biquadratic-V-art-#1}\centerline{\textbf{[#1.]}}\nobreak\smallskip}
\Article{1}
\[
\begin{array}{lll}
\text{Modulus}&M=A+Bi&AA+BB=D\\
\text{Rest}&m=a+bi,&aa+bb=d
\end{array}
\]
\[
\frac{mD}{M}=\mu=\alpha+\beta i=aA+bB+(Ab-Ba)i
\]
\[
\xi+\eta i=\pi,\quad\pi m=p=x+yi;\quad\pi M=P=X+Yi
\]
$\omega$ ist eine unbestimmte unendlich kleine reelle positive Größe.

\Article{2}
\begin{center}\emph{Vorbereitung.}\end{center}
I. Man sammle alle $\pi$, für die gilt:
\[
\begin{array}{l}
\xi\text{ nicht negativ und nicht größer als }\frac12\\
\eta\text{ positiv und kleiner als }\frac12\\
\text{Entweder }x\text{ oder }y\text{ eine ganze Zahl}\\
\text{Entweder }X\text{ oder }Y\text{ eine ganze Zahl}
\end{array}
\]
und bestimme für jedes $\pi$ die Größe $\varepsilon$ nach folgender Regel:
\[
\begin{array}{ll}
\text{Es sei }p^0&\text{die nächste durch }1+i\text{ teilbare ganze Zahl bei }p\\
P^0&\text{die nächste durch }1+i\text{ teilbare ganze Zahl bei }P
\end{array}
\]
und setze $\varepsilon=\pm1$, wobei das Zeichen stets dasselbe ist wie das Zeichen des imaginären Teils der Größe
\[
\frac{p-p^0}{P-P^0}(\alpha-\beta i)
\]

Es gelten die folgenden besonderen Regeln. Erste Klasse, $x$ und $X$ ganze Zahlen
\[
\begin{aligned}
\beta\xi&=-Bx+bX\\
\beta\eta&=-Ax+aX\\
\beta y&=-\alpha x+dX\\
\beta Y&=-Dx+\alpha X
\end{aligned}
\]

\SourceLeaf{358}
$\varepsilon=-1$, wenn $\beta$ positiv, $x$ gerade, $[y]$ gerade, $X$ gerade, $[Y]$ gerade ist oder wenn nur eine ungerade Anzahl dieser Bedingungen gilt.

$\varepsilon=+1$, wenn keine oder eine gerade Anzahl dieser Bedingungen gilt.

Zweite Klasse, $y$ und $Y$ ganze Zahlen
\[
\begin{aligned}
\beta\xi&=+Ay-aY\\
\beta\eta&=-By+bY\\
\beta x&=+\alpha y-dY\\
\beta X&=+Dy-\alpha Y
\end{aligned}
\]
$\varepsilon=-1$, wenn $\beta$ positiv, $[x]$ gerade, $y$ gerade, $[X]$ gerade, $Y$ gerade ist oder wenn eine ungerade Anzahl dieser Bedingungen gilt.

$\varepsilon=+1$, wenn keine oder eine gerade Anzahl gilt.

Dritte Klasse, $y$ und $X$ ganze Zahlen
\[
\begin{aligned}
\alpha\xi&=+By+aX\\
\alpha\eta&=+Ay-bX\\
\alpha x&=-\beta y+dX\\
\alpha Y&=+Dy-\beta X
\end{aligned}
\]
$\varepsilon=-1$, wenn $\alpha$ positiv ist und $[x],y,X,[Y]$ alle gerade sind oder wenn eine ungerade Anzahl dieser Bedingungen gilt.

$\varepsilon=+1$, wenn keine oder eine gerade Anzahl erfüllt ist.

Vierte Klasse, $x$ und $Y$ ganze Zahlen
\[
\begin{aligned}
\alpha\xi&=+Ax+bY\\
\alpha\eta&=-Bx+aY\\
\alpha y&=+\beta x+dY\\
\alpha X&=+Dx+\beta Y
\end{aligned}
\]
$\varepsilon=+1$, wenn $\alpha$ positiv ist und $x,[y],[X],Y$ alle gerade sind oder wenn eine ungerade Anzahl dieser Bedingungen gilt.

$\varepsilon=-1$, wenn keine oder eine gerade Anzahl davon gilt.

\SourceLeaf{359}
\Needspace{22\baselineskip}
\Article{3}
II. Man sammle alle $\pi$, für die gilt:
\[
\begin{array}{l}
\xi\text{ positiv und kleiner als }\frac12\\
\eta=\omega\\
\text{und entweder }X\text{ oder }Y\text{ eine ganze Zahl,}
\end{array}
\]
und setze $\varepsilon=\pm1$, wobei das Zeichen des imaginären Teils von
\[
\frac{M}{P-P^0}
\]
zu nehmen ist.

Besondere Regel: erste Klasse, $X$ ganz
\[
\begin{aligned}
A\xi&=+X+B\omega\\
Ax&=+aX-\beta\omega\\
Ay&=+bX-\alpha\omega\\
AY&=+BX+D\omega
\end{aligned}
\]
$\varepsilon=-1$, wenn $A$ positiv ist und $X$ und $[Y]$ gerade sind oder wenn nur eine Bedingung gilt.

$\varepsilon=+1$, wenn keine oder zwei gelten.

Zweite Klasse, $Y$ ganz
\[
\begin{aligned}
B\xi&=+Y-A\omega\\
Bx&=+aY-\alpha\omega\\
By&=+bY-\beta\omega\\
BX&=+AY-D\omega
\end{aligned}
\]
$\varepsilon=+1$, wenn $B$ positiv ist und $[X]$ und $Y$ gerade sind oder wenn nur eine Bedingung gilt.

$\varepsilon=-1$, wenn keine oder zwei gelten.

\clearpage
\Article{4}
III. Man sammle alle $\pi$, für die gilt:
\begin{quote}
$\xi$ und $\eta$ unterliegen denselben Bedingungen wie in II;\par
entweder $x$ oder $y$ ist ganz,
\end{quote}
und setze $\varepsilon=\pm1$ mit dem Zeichen des imaginären Teils von
\[
\frac{m}{p-p^0}
\]

\SourceLeaf{360}
Besondere Regeln: erste Klasse, $x$ ganz
\[
\begin{aligned}
a\xi&=+x+b\omega\\
ay&=+bx+d\omega\\
aX&=+Ax+\beta\omega\\
aY&=+Bx+\alpha\omega
\end{aligned}
\]
$\varepsilon=-1$, wenn $a$ positiv ist und $x,[y]$ beide gerade sind oder wenn nur eine Bedingung gilt.

$\varepsilon=+1$, wenn keine oder zwei gelten.

Zweite Klasse, $y$ ganz
\[
\begin{aligned}
b\xi&=+y-a\omega\\
bx&=+ay-d\omega\\
bX&=+Ay-\alpha\omega\\
bY&=+By+\beta\omega
\end{aligned}
\]
$\varepsilon=+1$, wenn $b$ positiv ist und $[x]$ und $y$ gerade sind oder wenn nur eine Bedingung gilt.

$\varepsilon=-1$, wenn keine gilt.

\clearpage
\Article{5}
IV. Man sammle alle $\pi$, für die gilt:
\[
\begin{array}{l}
\xi=\frac12+\frac12\omega\\
\eta\text{ positiv und kleiner als }\frac12\\
X\text{ oder }Y\text{ ganz,}
\end{array}
\]
und setze $\varepsilon=\pm1$ mit dem Zeichen des imaginären Teils von
\[
\frac{iM}{P-P^0}
\]
Besondere Regeln: erste Klasse, $X$ eine ganze Zahl,
\[
\begin{aligned}
2B\eta&=+A-2X+A\omega\\
2Bx&=-\beta+2bX-\beta\omega\\
2By&=+\alpha-2aX+\alpha\omega\\
2BY&=+D-2AX+D\omega
\end{aligned}
\]
$\varepsilon=+1$, wenn $B$ positiv ist und $X,[Y]$ gerade sind oder wenn eine Bedingung gilt,

$\varepsilon=-1$, wenn keine oder zwei Bedingungen gelten.

\SourceLeaf{361}
Zweite Klasse, $Y$ eine ganze Zahl
\[
\begin{aligned}
2A\eta&=-B+2Y-B\omega\\
2Ax&=+\alpha-2bY+\alpha\omega\\
2Ay&=+\beta+2aY+\beta\omega\\
2AX&=+D-2BY+D\omega
\end{aligned}
\]
\noindent $\varepsilon=+1$, wenn $A$ positiv ist und $[X],Y$ gerade sind oder wenn eine Bedingung gilt;\par
\noindent $\varepsilon=-1$, wenn keine oder zwei Bedingungen gelten.

\clearpage
\Article{6}
V. Man sammle alle $\pi$, für die gilt:
\[
\begin{gathered}
\xi,\eta\text{ denselben Bedingungen unterworfen sind wie in IV,}\\
\text{und wo }x\text{ oder }y\text{ eine ganze Zahl,}
\end{gathered}
\]
und setze $\varepsilon=\pm1$ mit dem Zeichen des imaginären Teils von
\[
\frac{im}{p-p^0}
\]

Besondere Regeln: erste Klasse, $x$ eine ganze Zahl
\[
\begin{aligned}
2b\eta&=+a-2x+a\omega\\
2by&=+d-2ax+d\omega\\
2bX&=+\beta+2Bx+\beta\omega\\
2bY&=+\alpha-2Ax+\alpha\omega
\end{aligned}
\]
\noindent $\varepsilon=+1$, wenn $b$ positiv ist und $x,[y]$ gerade sind oder wenn eine Bedingung gilt;\par
\noindent $\varepsilon=-1$, wenn keine oder zwei Bedingungen gelten.

Zweite Klasse, $y$ eine ganze Zahl
\[
\begin{aligned}
2a\eta&=-b+2y-b\omega\\
2ax&=+d-2by+d\omega\\
2aX&=+\alpha-2By+\alpha\omega\\
2aY&=-\beta+2Ay-\beta\omega
\end{aligned}
\]
\noindent $\varepsilon=+1$, wenn $a$ positiv ist und $[x],y$ gerade sind oder wenn eine Bedingung gilt;\par
\noindent $\varepsilon=-1$, wenn keine oder zwei Bedingungen gelten.
